%! Author = denis
%! Date = 20/08/2026

\begin{frame}{Como Descrever um Circuito Combinacional}

    Um mesmo circuito pode ser representado de diferentes formas:

    \[
        \boxed{\text{Tabela Verdade}}
        \Longleftrightarrow
        \boxed{\text{Expressão Booleana}}
        \Longleftrightarrow
        \boxed{\text{Circuito Lógico}}
    \]

    \vspace{0.5cm}

    Exemplo:

    \[
        F=A\cdot B
    \]

    \begin{center}
        \begin{tabular}{c c | c}
            \hline
            A & B & F \\
            \hline
            0 & 0 & 0 \\
            0 & 1 & 0 \\
            1 & 0 & 0 \\
            1 & 1 & 1 \\
            \hline
        \end{tabular}
    \end{center}

    Neste caso, o circuito é formado por uma porta AND.

\end{frame}
